\section{Cómo familiarizarse rápido con un proyecto de Agent}

\subsection{Síntesis de la metodología}

El autor sintetizó una ruta de tres pasos para comprender y familiarizarse rápido con un proyecto de Agent:

\begin{enumerate}
    \item \textbf{Comprender los conceptos básicos}: leer el artículo o la introducción al framework del proyecto, comprender el diseño del Workflow / Pipeline
    \item \textbf{Comprender el Prompt}: en un proyecto de Agent, gran cantidad de subprocesos son representados por el LLM; comprender el Prompt equivale a comprender la entrada, la salida y la lógica de procesamiento de cada subproceso
    \item \textbf{Usar herramientas de monitoreo}: usar herramientas como LangFuse para hacer trace de todas las llamadas a la API y comprender de forma visual el proceso completo
\end{enumerate}

\begin{knowledgebox}{El Prompt es la clave para comprender un proyecto de Agent}
En el Workflow de un Agent, muchos subprocesos son representados por el LLM. Si se entiende cada subproceso como $F(X; \text{Prompt})$ —donde $F$ es el LLM, $X$ es la Query y Prompt es el System/User Prompt—, comprender el Prompt equivale a comprender la lógica de procesamiento de cada subproceso.
\end{knowledgebox}

\subsection{Resumen de la sección}

Concepto $\rightarrow$ Prompt $\rightarrow$ herramienta de monitoreo: esta es una ruta eficiente para comprender un proyecto de Agent.

